\section{本章小结}
\label{sec:mcalc-summary}

本章开头的问题是：在流形上怎样求导、怎样积分，内部的变化怎样汇总到边界。第一部分只用光滑结构回答。微分形式是对多个切向量的交替读数，楔积按次数满足交换符号规则，拉回沿映射反向作用且对任何光滑映射都有定义，外微分把$k$形式变为$k+1$形式，对函数就是$\mathrm df$，与楔积、拉回相容，并满足$\mathrm d\circ\mathrm d=0$。闭而不恰当的形式，如穿孔平面上的$(x_1\,\mathrm dx_2-x_2\,\mathrm dx_1)/r^2$，反映流形的拓扑。

积分最高次形式还需要定向：形式的系数按带符号的Jacobian行列式变换，只有在定向图册中它才与积分换元公式中的绝对值一致。球面的两张球极投影图要调整一个坐标的符号才构成定向图册，M\"obius带则根本没有定向。边界的诱导定向采用“外法向在前”的约定，在这一约定下，Stokes公式把$\mathrm d\omega$在流形上的积分与$\omega$在边界上的积分联系起来，对紧支撑的$(n-1)$形式成立，紧支撑与定向都不可省略；Green、Gauss与经典Stokes公式是它在平面区域、三维区域与有向曲面上的特例。

第二部分加入黎曼度量。体积密度$\sqrt{\det\bm G}\,\mathrm d\bm x$以Jacobian的绝对值变换，因而不需要定向，在可定向时对应体积形式。由它得到的函数积分是一个测度积分，与定向无关，随度量改变，也不同于样本的采样测度。梯度用度量把$\mathrm df$变成向量，Hessian用Levi-Civita联络扣除方向场的变化，是对称双线性型；散度是协变导数的迹，坐标式只含体积密度；Laplace--Beltrami算子是梯度的散度，在本书约定下非正。球面高度函数把这些运算串在一起：$\operatorname{grad}h$沿经线指向北极，$\operatorname{Hess}h=-hg$，$\Delta h=-2h$，三个结果彼此印证。

梯度、散度与散度定理合成分部积分公式，边界项在闭流形、紧支撑或Dirichlet、Neumann条件下消失，非紧时须另行核对；它也给出弱梯度与场方程弱形式的定义。在紧流形上，有限坐标覆盖、单位分解和黎曼体积把$L^p$与Sobolev空间搬了过来，空间本身不依赖度量的选择，完备性与稠密性得以保留。非紧时须区分局部与整体的控制，欧氏的紧嵌入等结论不能无条件继承。这些运算为后续的流形优化与场方程提供了合法使用的条件。
